\section{贡献二：Dialogue-Aware Search 与 piKL}

有了可控语言，还需要一个会与人类共处的 planner。纯 self-play 可以追求高收益，却可能收敛到人类从未见过的约定；纯 imitation 容易被理解，却只复制平均人类错误，并可能被一条诱导消息操纵。piKL 把这两个极端放进同一个优化目标。

\subsection{从贡献列表进入规划问题}

slide 31 将第二项贡献高亮：planner 必须同时考虑 dialogue 与 human behavior。它不能把消息当作无关文本，因为可信的承诺会改变他人行动；也不能无条件服从消息，因为对手可能建议一个对 Cicero 不利的计划。

\lecturefigure{slide-31-main-contribution-planning.jpg}{Main contribution：面向对话与人类行为的 planning}{官方录像恢复幻灯片 V031；Stanford CS25 Lecture 14}

\subsection{Search 与 imitation 的互补失败}

slide 32 左侧的 pure search / RL 追求 return，却不受人类惯例约束；右侧 imitation 贴近数据，却继承弱行动，并把语言相关性学成表面相关。中间的 piKL 像一条“腌黄瓜”连接两端：以 human policy 为 anchor，再允许 search 在价值驱动下偏离。

\lecturefigure{slide-32-dialogue-aware-search-rl.jpg}{Dialogue-aware search：在纯优化与纯模仿之间寻找可用策略}{官方录像恢复幻灯片 V032；Stanford CS25 Lecture 14}

\begin{warningbox}{两个极端各自会怎样坏掉}
纯 self-play policy 可能采用只有自己复制体理解的 convention，面对人类时不可协调；纯 imitation policy 则可能忠实复现人类次优行为，并对“请把你的关键领地让给我”这类消息反应过度。可沟通 agent 既要让人类可预测，也要保留拒绝坏建议的能力。
\end{warningbox}

\subsection{piKL：价值与人类先验的显式权衡}

令 $\pi_H(a\mid s,h,d)$ 为从人类数据学习的 anchor policy，$Q^{\pi}(s,a)$ 为 search 下动作价值，piKL 的单状态直觉目标可写成
\begin{equation}
\pi^*
=\arg\max_{\pi}
\left[
\mathbb{E}_{a\sim\pi}Q^{\pi}(s,a)
-\lambda D_{\mathrm{KL}}\!\left(\pi(\cdot\mid s)\,\|\,\pi_H(\cdot\mid s)\right)
\right].
\end{equation}
$\lambda$ 控制偏离人类行为的成本。$\lambda$ 太小，策略可能强但怪异；$\lambda$ 太大，系统接近原始 imitation，难以纠正数据中的次优行动。它不是把 human-likeness 当作最终目标，而是把人类 convention 作为规划 prior 和他人预期模型。

\lecturefigure{slide-33-pikl.jpg}{piKL：以 human imitation policy 为锚的 KL-regularized search}{官方录像恢复幻灯片 V033；Jacob et al., ICML 2022}

\begin{knowledgebox}{读图：三类模型在同一个循环中的位置}
图左的 human policy 提供 anchor，中央 search / RL 计算价值改善，右侧环境反馈更新状态。KL 正则不是简单 post-processing，而是在每次策略更新中持续限制分布漂移。结果应当同时评估 exploitability / return 与 human predictability，不能只报一边。
\end{knowledgebox}

\teachervoice{讲者指出，类似 KL penalty 并不新鲜；Cicero 更特别的地方是把 anchor policy 同时当成“人类会怎么做”和“人类认为别人会怎么做”的模型。这里出现了 common knowledge / theory-of-mind 味道：一个行动的价值取决于别人是否把它视为合理 convention。}

\subsection{对话怎样真实改变策略}

slide 34 连续展示三种对话后的策略变化。第一种消息表示合作方要退出某方向，Cicero 因而回防；第二种表达不信任，使它降低依赖该合作的行动；第三种提出对 Cicero 不利的建议，value model 阻止系统盲从。关键不是“每条消息都改变动作”，而是只有可信且有战略后果的消息才应改变概率分布。

\lecturefigure{slide-34-strategic-response-sequence.jpg}{Strategic response sequence：对话信号如何改变或不改变行动 policy}{官方录像恢复幻灯片 V034；Stanford CS25 Lecture 14}

\begin{importantbox}{Dialogue-aware 不等于 instruction-following}
对手消息是带策略目的的 observation，不是系统指令。planner 应把它作为更新 belief 的证据，再用自身 value function 决定是否响应。这个区分使 Cicero 能合作，也能拒绝 manipulation。
\end{importantbox}

\subsection{Modeling Human Policies}

上一节已经展示 policy 会随可信对话移动；要实现这种更新，系统必须先把文本转成行为预测。本节说明 planner 如何消费语言，而不只是如何生成语言。slide 35 把自然语言历史输入 BART，再预测人类 orders。它体现了一个常被忽略的接口：语言模型不只用于生成消息，也可把对话转化为对其他玩家行为的概率预测。预测分布随后进入 search，形成 dialogue 到 action 的可审计路径。

\lecturefigure{slide-35-model-human-policies.jpg}{Modeling human policies：从对话与状态预测人类行动分布}{官方录像恢复幻灯片 V035；Stanford CS25 Lecture 14}

\teachervoice{老师用多个实例说明，Cicero 并不会因为对方“说了”就改变计划。只有消息与局势、历史和对方利益共同支持某种行为时，policy 才明显移动；value estimate 则负责挡住容易利用 imitation model 的诱导。}

\subsection{本章小结}

piKL 通过 KL 正则，把 search 的战略强度与人类惯例的可协调性结合起来。对话感知规划的本质也不是服从语言，而是把语言转化为他人行动分布，再在价值约束下更新自己的策略。

